\documentclass[11pt]{article}

\usepackage[margin=1in]{geometry}
\usepackage{amsmath}
\usepackage{array}
\usepackage{booktabs}
\usepackage{longtable}
\usepackage[T1]{fontenc}
\usepackage{lmodern}
\usepackage[hidelinks]{hyperref}

\title{Lecture Notes on the IEEE 9-Bus AC Power-Flow Problem}
\author{Generated from the loadflow-lite MATPOWER fixtures}
\date{\today}

\newcommand{\pu}{\mathrm{p.u.}}
\newcommand{\MW}{\mathrm{MW}}
\newcommand{\MVAr}{\mathrm{MVAr}}

\begin{document}

\maketitle

\begin{abstract}
The IEEE 9-bus system is a compact benchmark for the steady-state AC
power-flow problem. It is small enough to inspect by hand, but it still
contains the main modeling features that make power flow nontrivial:
multiple generator buses, voltage magnitude control, nonlinear active-
and reactive-power balance equations, branch losses, and reactive-power
exchange with transmission lines. These notes formulate the problem,
describe the specific 9-bus network used in this repository, and discuss
the solved operating point recorded in the checked-in MATPOWER-derived
fixtures.
\end{abstract}

\section{Context and Relevance}

Power-flow analysis is the basic steady-state calculation behind many
power-system engineering tasks. Given a network model and specified
generation, load, and voltage-control data, the AC power-flow problem
computes the complex voltage at each bus and the real and reactive power
flows on each branch. The result is not a dynamic trajectory; it is a
snapshot of an operating point under balanced, sinusoidal, three-phase
conditions.

This calculation matters because many operational and planning decisions
depend on quantities that are not specified directly in the input data:
voltage magnitudes at load buses, bus voltage angles, slack-bus real-power
output, generator reactive-power output, branch flows, and system losses.
A candidate dispatch may satisfy a real-power energy balance in aggregate
while still producing unacceptable voltage profiles, excessive line flows,
or reactive-power requirements. AC power flow is therefore the entry point
for contingency analysis, voltage-security assessment, optimal power flow,
state estimation validation, and many simulation studies.

The IEEE 9-bus case is useful pedagogically because it is small but not
degenerate. It has one slack bus, two PV generator buses, six PQ buses, and
nine branches on a 100 MVA base. The topology contains a generator/load
network rather than a single radial feeder, and the solution exhibits both
real-power losses and significant reactive-power effects from line charging.

\section{Mathematical Formulation}

\subsection{Network Model}

Let the system have $n$ buses and complex bus-voltage phasors
\[
  V_i = |V_i| e^{j\theta_i}, \qquad i = 1,\ldots,n.
\]
The nodal admittance matrix is
\[
  Y_{\mathrm{bus}} = G + jB,
\]
where $G_{ik}$ and $B_{ik}$ are the conductance and susceptance entries.
At each bus, the net complex power injection is
\[
  S_i = P_i + jQ_i = V_i I_i^*.
\]
Using $I = Y_{\mathrm{bus}}V$, the standard polar-form AC power-flow
equations are
\begin{align}
  P_i &=
  |V_i| \sum_{k=1}^n |V_k|
  \left[
    G_{ik}\cos(\theta_i-\theta_k)
    + B_{ik}\sin(\theta_i-\theta_k)
  \right], \\
  Q_i &=
  |V_i| \sum_{k=1}^n |V_k|
  \left[
    G_{ik}\sin(\theta_i-\theta_k)
    - B_{ik}\cos(\theta_i-\theta_k)
  \right].
\end{align}

The specified net injections are formed from positive generation and
positive load:
\[
  P_i^{\mathrm{spec}} = P_{G,i} - P_{D,i}, \qquad
  Q_i^{\mathrm{spec}} = Q_{G,i} - Q_{D,i}.
\]
All real and reactive powers in these notes are reported on the 100 MVA
base used by the fixture files unless explicitly shown in MW or MVAr.

\subsection{Bus Types and Unknowns}

The power-flow equations are solved with different specifications at
different bus types:
\begin{itemize}
  \item \textbf{Slack bus}: specifies $|V|$ and $\theta$. Its real and
  reactive generation are outputs of the solution. It absorbs the
  imbalance caused by losses and numerical consistency.
  \item \textbf{PV bus}: specifies real-power generation $P_G$ and voltage
  magnitude $|V|$. The voltage angle $\theta$ and reactive generation
  $Q_G$ are solved.
  \item \textbf{PQ bus}: specifies real and reactive load, and any fixed
  generation if present. Both $|V|$ and $\theta$ are solved.
\end{itemize}

For a Newton-Raphson solution, the nonlinear equations are written as
mismatches. For non-slack buses,
\[
  \Delta P_i = P_i^{\mathrm{spec}} - P_i(V,\theta),
\]
and for PQ buses,
\[
  \Delta Q_i = Q_i^{\mathrm{spec}} - Q_i(V,\theta).
\]
The Jacobian maps small changes in voltage angles and PQ-bus voltage
magnitudes to these active- and reactive-power mismatches. Iteration
continues until the mismatches are sufficiently small.

\section{The Specific IEEE 9-Bus Network}

The fixture used here is \texttt{data/case9.json}, with solved values from
\texttt{data/case9\_solution.json}. The case has three generator buses:
bus 1 is the slack bus, while buses 2 and 3 are PV buses. Loads are located
at buses 5, 7, and 9. The total load is
\[
  P_D = 315.00~\MW, \qquad Q_D = 115.00~\MVAr.
\]
The case-side scheduled generation sums to $320.30~\MW$ and
$22.62~\MVAr$, but the solved generator outputs differ because the slack
bus and PV-bus reactive outputs are determined by the AC solution.

\begin{table}[htbp]
\centering
\caption{Bus data and solved voltages for the IEEE 9-bus case.}
\begin{tabular}{@{}rrrrrrrrr@{}}
\toprule
Bus & Type & $P_D$ & $Q_D$ & $P_G^{\mathrm{case}}$ & $Q_G^{\mathrm{case}}$ &
$|V|$ & $\theta$ & $P_G^{\mathrm{sol}}$ \\
 & & MW & MVAr & MW & MVAr & p.u. & deg & MW \\
\midrule
1 & slack & 0.00 & 0.00 & 72.30 & 27.03 & 1.0400 & 0.0000 & 71.64 \\
2 & PV & 0.00 & 0.00 & 163.00 & 6.54 & 1.0250 & 9.2800 & 163.00 \\
3 & PV & 0.00 & 0.00 & 85.00 & -10.95 & 1.0250 & 4.6648 & 85.00 \\
4 & PQ & 0.00 & 0.00 & 0.00 & 0.00 & 1.0258 & -2.2168 & 0.00 \\
5 & PQ & 90.00 & 30.00 & 0.00 & 0.00 & 1.0127 & -3.6874 & 0.00 \\
6 & PQ & 0.00 & 0.00 & 0.00 & 0.00 & 1.0324 & 1.9667 & 0.00 \\
7 & PQ & 100.00 & 35.00 & 0.00 & 0.00 & 1.0159 & 0.7275 & 0.00 \\
8 & PQ & 0.00 & 0.00 & 0.00 & 0.00 & 1.0258 & 3.7197 & 0.00 \\
9 & PQ & 125.00 & 50.00 & 0.00 & 0.00 & 0.9956 & -3.9888 & 0.00 \\
\bottomrule
\end{tabular}
\end{table}

\begin{table}[htbp]
\centering
\caption{Branch parameters for the 9-bus case. All series and charging
parameters are in p.u. on the system base.}
\begin{tabular}{@{}rrrrrr@{}}
\toprule
Branch & From & To & $r$ & $x$ & $b$ \\
\midrule
1 & 1 & 4 & 0.00000 & 0.05760 & 0.00000 \\
2 & 4 & 5 & 0.01700 & 0.09200 & 0.15800 \\
3 & 5 & 6 & 0.03900 & 0.17000 & 0.35800 \\
4 & 3 & 6 & 0.00000 & 0.05860 & 0.00000 \\
5 & 6 & 7 & 0.01190 & 0.10080 & 0.20900 \\
6 & 7 & 8 & 0.00850 & 0.07200 & 0.14900 \\
7 & 8 & 2 & 0.00000 & 0.06250 & 0.00000 \\
8 & 8 & 9 & 0.03200 & 0.16100 & 0.30600 \\
9 & 9 & 4 & 0.01000 & 0.08500 & 0.17600 \\
\bottomrule
\end{tabular}
\end{table}

Topologically, buses 1, 2, and 3 connect generator sources into a meshed
transmission network through buses 4, 8, and 6 respectively. The load buses
5, 7, and 9 are embedded in the loop
\[
  4 \rightarrow 5 \rightarrow 6 \rightarrow 7 \rightarrow 8
  \rightarrow 9 \rightarrow 4.
\]
This loop is important: real-power injections do not simply flow along a
single path. They divide according to the network impedances and the bus
voltage-angle profile.

\section{Discussion of the Solved Operating Point}

The solved generator outputs are
\[
\begin{aligned}
  P_G^{\mathrm{sol}} &= 319.641~\MW, \\
  Q_G^{\mathrm{sol}} &= 22.840~\MVAr.
\end{aligned}
\]
Since the total real load is $315.000~\MW$, the real-power loss is
\[
  P_{\mathrm{loss}} = 319.641 - 315.000 = 4.641~\MW.
\]
The slack bus supplies $71.641~\MW$, slightly less than the case-side
scheduled value of $72.300~\MW$, because the PV generator dispatch and the
network losses determine the final balance. Buses 2 and 3 hold their
scheduled real-power outputs of $163.000~\MW$ and $85.000~\MW$.

The voltage magnitudes remain close to nominal. The highest solved voltage
is the slack-bus voltage, $1.0400~\pu$, while the lowest voltage is at bus
9, $0.9956~\pu$. The lowest-voltage bus is also the largest load bus
($125~\MW$ and $50~\MVAr$), which is consistent with the expected voltage
drop around the transmission loop. The voltage angles range from
$-3.9888^\circ$ at bus 9 to $9.2800^\circ$ at bus 2. In an AC transmission
network, active-power transfer is strongly tied to voltage-angle
differences, so these angles are not merely reporting variables; they are
part of the solved state that explains the branch real-power flows.

\begin{longtable}{@{}rrrrrrrr@{}}
\caption{Solved branch powers. The ``from'' and ``to'' quantities follow
the fixture branch orientation; negative values mean flow enters that end
of the branch.}\label{tab:branch-flows}\\
\toprule
From & To & $P_f$ & $Q_f$ & $P_t$ & $Q_t$ & $P_{\ell}$ & $Q_{\ell}$ \\
 & & MW & MVAr & MW & MVAr & MW & MVAr \\
\midrule
\endfirsthead
\toprule
From & To & $P_f$ & $Q_f$ & $P_t$ & $Q_t$ & $P_{\ell}$ & $Q_{\ell}$ \\
 & & MW & MVAr & MW & MVAr & MW & MVAr \\
\midrule
\endhead
1 & 4 & 71.641 & 27.046 & -71.641 & -23.923 & 0.000 & 3.123 \\
4 & 5 & 30.704 & 1.030 & -30.537 & -16.543 & 0.166 & -15.513 \\
5 & 6 & -59.463 & -13.457 & 60.817 & -18.075 & 1.354 & -31.531 \\
3 & 6 & 85.000 & -10.860 & -85.000 & 14.955 & 0.000 & 4.096 \\
6 & 7 & 24.183 & 3.120 & -24.095 & -24.296 & 0.088 & -21.176 \\
7 & 8 & -75.905 & -10.704 & 76.380 & -0.797 & 0.475 & -11.502 \\
8 & 2 & -163.000 & 9.178 & 163.000 & 6.654 & 0.000 & 15.832 \\
8 & 9 & 86.620 & -8.381 & -84.320 & -11.313 & 2.300 & -19.694 \\
9 & 4 & -40.680 & -38.687 & 40.937 & 22.893 & 0.258 & -15.794 \\
\bottomrule
\end{longtable}

Table~\ref{tab:branch-flows} also illustrates the importance of sign
conventions. For example, on branch $8$--$2$, the fixture orientation is
from bus 8 to bus 2, but $P_f=-163.000~\MW$ and $P_t=163.000~\MW$. The
real power therefore flows from bus 2 toward bus 8, not from bus 8 toward
bus 2. Similarly, on branch $7$--$8$, the negative $P_f$ indicates that
real power enters the branch at bus 8 and leaves at bus 7.

Summing $Q_f+Q_t$ over all branches gives approximately
$-92.160~\MVAr$. This is not a violation of reactive-power conservation;
it reflects the shunt charging susceptance in the line models. Several
branches inject more reactive power through their charging elements than
they absorb in series reactive drop, so the net branch reactive ``loss''
can be negative. This is one of the reasons AC power flow must model
reactive power explicitly rather than treating it as a small correction to
a real-power-only calculation.

\section{Takeaways}

\begin{enumerate}
  \item The IEEE 9-bus problem is a minimal but meaningful AC power-flow
  benchmark: it includes slack, PV, and PQ bus behavior in a small meshed
  network.
  \item The unknown state consists of voltage angles at non-slack buses and
  voltage magnitudes at PQ buses. Generator reactive powers and slack-bus
  power are outputs, not fixed inputs.
  \item The solved real-power generation exceeds total real load by
  $4.641~\MW$, which is the network real-power loss at this operating
  point.
  \item The branch-flow signs must be interpreted relative to the branch
  orientation. Negative ``from'' power means the actual real-power flow is
  opposite to the stored from-to direction.
  \item Reactive-power behavior is qualitatively different from real-power
  loss accounting because line charging can inject reactive power.
\end{enumerate}

\end{document}
